%% 06-discussion.tex --- included by ../sfc.tex
\section{Discussion and Limitations}
\label{sec:discussion}

Section~\ref{sec:results} reports several results that are individually
clear but that read as tensions when placed side by side. We resolve them
here, then state what the evaluation does not establish.

\textbf{What the plan abstraction buys, and what commitment buys.}
These are separable and the paper's own terminology has to keep them
apart. The \emph{plan-oriented} design --- generating whole inter-core
actions as candidates and ranking them on gain with preparation cost
demoted to a tie-break --- is what distinguishes \PORTER{}'s remote stage,
and it is the configuration that produces the EPR reductions of
Section~\ref{sec:rq1}. \emph{Commitment} is the narrower rule that the
selected plan is not reopened until it is ready. Only the latter is
ablated in Table~\ref{tab:ablation}: the hop-wise column re-derives and
re-scores the same plan set every iteration, so it changes when the choice
is revisited and nothing else. That it reaches the same EPR on the same
number of routing moves says commitment is not where the quality comes
from; it does not say the plan abstraction is idle, because no column
removes it. Isolating the abstraction would need a hop-granular candidate
set, which is a different router and which we have not built. This is the
clearest gap in our attribution.

\textbf{Layout below about a hundred physical qubits, horizon above it.}
\PORTER{} trails \dSABRE{} under RQ2 at $\le 108$ physical qubits, and
controlling the layout while matching the pass budget erases the
difference (Sections~\ref{sec:shared_results} and~\ref{sec:heavyhex}).
On the scale-up series neither the layout control nor extra refinement
passes help --- \dSABRE{}'s layout is worse for \PORTER{} there than its
own --- while selecting a longer lookahead horizon moves \PORTER{} from
$1.45\times$ behind to $0.72\times$ ahead. The two regimes are consistent
with the horizon being an absolute gate count that does not grow with the
core count $K$, so that a plan at twenty cores is scored over
proportionally less of the circuit than at four. We have not tested a
$K$-scaled horizon, so that remains a hypothesis rather than a finding.

\textbf{Why \TeleData{}-only can beat the full action set.} Packets
clearly beat the single telegates they degrade into, yet removing gate
teleportation entirely is cheaper still on the packet-heaviest circuits
at 80 qubits. We read this as an interaction between the action
vocabulary and the scorer rather than as a failure of packetisation. The
two are scored in different currencies: a packet's gain is a count of
remote CX gates it resolves immediately, whereas a \TeleData{} plan's is a
sum of core-distance reductions, most of it discounted over a bounded
horizon. Nothing establishes that one resolved gate and one unit of
distance reduction are commensurate, and where a teleport's repayment
extends past the horizon the comparator cannot see it. The ablation is
evidence that this mismatch costs EPR pairs at scale; it is not evidence
that a pair spread over a run of gates is a bad idea.

\textbf{What the completion record does and does not show.}
Section~\ref{sec:sfc-progress} separates correctness and termination,
which \PORTER{} guarantees by construction, from completion, which it does
not. Every instance in this evaluation completes under the selected-horizon
protocol, but that is a property of the protocol and the benchmark rather
than of the router: at a fixed horizon, heavy-hex star Multiplier
livelocks, and the extra candidate is a way of not meeting the failure
rather than of excluding it. Two of the four failure modes were observed
here --- livelock at a fixed horizon, and infeasibility on an external
pytket-dqc layout that fills a core --- while capacity deadlock was
observed only with relief disabled. \dSABRE{}'s safe mode carries a
termination-with-success proof under stated architectural hypotheses, and
supplying a comparable guarantee is the clearest piece of theory this
design leaves open.

\subsection{Threats to Validity}
\label{sec:threats}

\textbf{Benchmark coverage.} The main suites are nine MQT-Bench families
at three widths; the scale-up series is two families (QFT and QPEexact)
only, so nothing here supports a claim about scaling behaviour on
irregular or randomly structured circuits at those widths. The two
heavy-hex networks are two core graphs, not a survey of architectures.

\textbf{Parameter rules fitted in-sample.} The horizon candidate sets and
the density threshold were chosen by inspecting these same instances
(Section~\ref{sec:horizon}). The magnitude of the horizon-selection gain
is therefore in-sample. The portfolio construction bounds the downside ---
the configured horizon is always a candidate and wins ties --- but an
out-of-sample validation has not been performed. The same caution applies
in weaker form to $f_\mathrm{min}$, $m_\mathrm{assist}$ and the packet
window, which were fixed early and not swept.

\textbf{Reconstructed benchmark instances.} The 80-qubit AE and QNN
circuits are not produced by the stock generator and were reconstructed
(\appSuite{}). We release the construction scripts and the files, and
validate the QNN reconstruction by exact CNOT-sequence match against the
64-qubit file, but they remain our circuits rather than the generator's.

\textbf{Patched baseline binary.} All \TeleSABRE{} numbers use a binary
patched to avoid an out-of-range index in its initial-layout pass
(Section~\ref{sec:baselines}). The patch changes placement only for qubits
the stock binary leaves unplaced, and architecture-H results, where the
stock binary does not crash, are unchanged by it; we have not otherwise
audited its effect.

\textbf{An evolving concurrent baseline.} \dSABRE{} was revised during
this work. We pin an exact revision and re-run rather than quoting, but a
later revision may move the comparison again.

\textbf{Unequal native budgets.} The own-protocol comparison gives
\dSABRE{} nine to eighteen routes against \PORTER{}'s three. This is each
compiler's own convention rather than a controlled budget, which is why
RQ2 exists and why compile time is reported everywhere.

\textbf{Aggregation over failures.} Geometric means are taken over jointly
completed instances, so a baseline that fails hard instances is scored
only on easier ones. Completion counts are reported alongside every such
mean, but the two must be read together.

\textbf{Implementation and timing.} \PORTER{} is Python and the baselines
are Python and C respectively, so per-pass times compare implementations
as much as algorithms. Compile times are wall-clock on one machine.
Horizon selection additionally guards its speculative routes with a
wall-clock budget, so on a loaded machine it may abandon a candidate that
would otherwise have won: the protocol's output is reproducible only up to
that budget, though each individual route is exactly reproducible.

\textbf{Cost model.} One EPR pair per inter-core action, independent of
hop count and generation protocol; link contention, entanglement latency
and fidelity, purification, and classical communication cost are all
outside the model (Section~\ref{sec:scope}). Conclusions are relative to
that model.
